\begin{definition}
\label{definition-galois-category}
\begin{reference}
Different from the definition in \cite[Expos\'e V, Definition 5.1]{SGA1}.
Compare with \cite[Definition 7.2.1]{BS}.
\end{reference}
Let $\mathcal{C}$ be a category and let $F : \mathcal{C} \to \textit{Sets}$
be a functor. The pair $(\mathcal{C}, F)$ is a {\it Galois category} if
\begin{enumerate}
\item $\mathcal{C}$ has finite limits and finite colimits,
\item
\label{item-connected-components}
every object of $\mathcal{C}$ is a finite (possibly empty)
coproduct of connected objects,
\item $F(X)$ is finite for all $X \in \Ob(\mathcal{C})$, and
\item $F$ reflects isomorphisms\footnote{Namely, given a morphism
$f$ of $\mathcal{C}$ if $F(f)$ is an isomorphism, then
$f$ is an isomorphism.} and is exact\footnote{This means that
$F$ commutes with finite limits and colimits, see
Categories, Section \ref{categories-section-exact-functor}.}.
\end{enumerate}
Here we say $X \in \Ob(\mathcal{C})$ is connected if
it is not initial and for any monomorphism $Y \to X$
either $Y$ is initial or $Y \to X$ is an isomorphism.
\end{definition}